% Folder 83, batch 10 (archivists' pages 181-200).
% Pass: Opus 5.5 (claude-opus-5-5), 2026-10-10 — first pass, unchecked against the pages by a human.
%
% Page 181 is blank; pages 195 and 197 are typescript pages (L 19, L 17) of a
% non-mathematical text by Grothendieck with no mathematics on them, skipped.
% Pages 194, 196, 198-200 are written on the versos of such typescripts; only
% his mathematics is transcribed. Subject: "(I bis)", March-July 1986 — the
% multiplication of a rank-4 Azumaya algebra recovered from det, 1 and the
% orientation e: tau = Tr xyz and tau' = Tr zyx satisfy tau+tau' = Sigma,
% tau tau' = Pi (explicit polynomials in b, c, t), tau - tau' = eps_0 Delta with
% Delta^2 a Gram determinant; the ring of G-invariant functions of three
% variables as a quadratic algebra over k[b,c,t] with S_3 action; then the
% four-variable case and the S_4 symmetries of the t_ijk (cube pictures).

\documentclass[11pt,a4paper]{article}
\input{../preamble/grothendieck}

\folder{83}
\batch{10}
\pages{181}{200}
\foldertitle{[Icosaèdres, bi-icosaèdres et algèbres d'Azumaya de rang 4] : notes manuscrites (1982-1983, 1986, s.d.).}
\dating{1982-1986}
\watermark{Édition de démonstration}

\begin{document}

\section{(I bis)}
\note{titre cerclé de sa main en tête de la p.~182, qui ouvre une nouvelle suite.}

\page{182}
\marginal{\uncertain{Mars} 1986 --- Juillet 1986}

$\mathcal{A}$ alg.\ d'Az.\ de rang 4 sur $k$, on veut expliciter sa structure
multiplicative en termes de la forme $\det$, de l'élément $1$, et de la
structure \emph{spéciale} orthogonale $e \in \hat{\Lambda}^4 \mathcal{A}$
\struck{\ill{}} \add{\ill{}} associée à $\det$ ($e \otimes e = \mathrm{dis}(\det)$).
Comme on connaît \add{$\det$ et $1$,} la forme $\mathrm{Tr}\, xy$ par la formule
\[
  \det(x+y) = \det x + \det y + \mathrm{Tr}\, x\, \mathrm{Tr}\, y - \mathrm{Tr}\, xy ,
\]
\[
  \varphi_{\det}(x,y) = \mathrm{Tr}\, x\, \mathrm{Tr}\, y - \underbrace{\mathrm{Tr}\, xy}_{\varphi(x,y)} ,
\]
\struck{pour} et que celle-ci est non dég., pour connaître $xy$, il suffit de connaître
\[
  \mathrm{Tr}\, xyz \qquad \forall z \in \mathcal{A} .
\]
\marginal{NB $\mathrm{Tr}$ \ill{} en termes de $\det$ : $\mathrm{Tr}(x) = \varphi_{\det}(x,1)$, \ill{} i.e.}
\note{le premier indice de $\varphi$ est écrit « det » sur un mot biffé illisible.}

Donc je cherche une formule pour $\mathrm{Tr}\, xyz$. Comme $\mathrm{Tr}(xy)z$ et
$\mathrm{Tr}(yx)z$ ne sont pas égaux en général, i.e.\ la structure mult.\
\add{de $\mathcal{A}$} n'est pas commutative \struck{(alors} on sent que $\det$, $1$ et
$e$ ne \uncertain{suffiront} à expliciter $\mathrm{Tr}\, xyz$ car ils ne déterminent la
structure multiplicative que mod passage à l'inverse, i.e.\ mod l'involution
\[
  x \longmapsto x' = \mathrm{Tr}(x).1 - x .
\]

On voit aussi que $\mathrm{Tr}\, xyz$ ne dépend que du triple $\{x,y,z\}$ et d'un
ordre circulaire sur l'ens.\ des trois termes. On a, \struck{pour} si
\[
  \begin{aligned}
    \tau &= \mathrm{Tr}\, xyz \quad (= \mathrm{Tr}\, yzx = \mathrm{Tr}\, zxy) \\
    \tau' &= \mathrm{Tr}(zyx) \quad (= \mathrm{Tr}(yxz) = \mathrm{Tr}(xzy))
  \end{aligned}
\]

\page{183}
la formule
\[
  \tau + \tau' = \mathrm{Tr}\,(xy+yx)z
\]
or
\[
  xy + yx = \underbrace{\mathrm{Tr}(y)}_{b'} x + \underbrace{\mathrm{Tr}(x)}_{b} y
  + \bigl(\underbrace{\mathrm{Tr}(xy)}_{t''} - \underbrace{\mathrm{Tr}\, x}_{b}\, \underbrace{\mathrm{Tr}\, y}_{b'}\bigr) 1
\]
\marginal{$\mathrm{Tr}\, x = b$, $\mathrm{Tr}\, y = b'$, $\mathrm{Tr}\, z = b''$ ;
$\det x = c$, $\det y = c'$, $\det z = c''$ ;
$\mathrm{Tr}\, yz = t$, $\mathrm{Tr}\, zx = t'$, $\mathrm{Tr}\, xy = t''$}
donc
\[
  \text{\struck{$(xy)z+$}}\quad (xy+yx)z = b'xz + byz + (t''-bb')z
\]
d'où
\[
  \boxed{\tau + \tau' = bt + b't' + b''t'' - bb'b''}
  \overset{\text{déf}}{=} \Sigma(b,b',b'',t,t',t'')
\]
\note{les deux premiers termes de l'encadré sont surchargés ; la lecture $bt + b't'$ est celle que donne le calcul de la ligne précédente.}

On devine que $\tau\tau'$ sera également un polynôme en
$b,b',b'',c,c',c'',t,t',t''$, de sorte que $\tau,\tau'$ s'explicitent (ici par
\uncertain{permutation} \struck{près}) comme racines d'un polynôme \add{unitaire} du second degré,
dont les coeff.\ $B$ ($= \tau+\tau'$) et $C$ ($= \tau\tau'$) sont des polynômes en les
variables $b^{(i)}, c^{(i)}, t^{(i)}$.

Mais plutôt que de chercher une expression pour $\tau\tau'$, je vais essayer de les
\add{[$\tau$ et $\tau'$]} déterminer directement en termes de $\det$, $1$ et de $e$.
Je suis guidé par le fait que $xy$ et $yx$ sont \uncertain{symétriques} mod $1, x, y$
--- donc si \emph{$z$ est orth.\ à $1, x, y$} \add{(pour $\varphi$)} \struck{$\det$} on aura
\[
  \mathrm{Tr}(xyz) = - \mathrm{Tr}(yxz) .
\]
\note{un long trait oblique traverse la moitié inférieure de la page sans annuler le texte, que la suite utilise.}
Pour se ramener à ce cas, on construit $\alpha, \beta, \gamma \in k$ tels que
\[
  z' = z - \alpha 1 - \beta x - \gamma y
\]
soit orth.\ à $1, x, y$, i.e.\ \note{dans la troisième ligne, avant $b'\alpha$, une première écriture biffée illisible.}
\[
  \begin{aligned}
    &b'' - 2\alpha - b\beta - b'\gamma = 0 \\
    &t' - b\alpha + (2c - b^2)\beta - t''\gamma = 0 \\
    &t - b'\alpha - t''\beta + (2c' - b'^2)\gamma = 0
  \end{aligned}
\]
i.e.
\[
  \left\{
  \begin{aligned}
    2\alpha + b\beta + b'\gamma &= b'' \\
    b\alpha + (b^2 - 2c)\beta + t''\gamma &= t' \\
    b'\alpha + t''\beta + (b'^2 - 2c')\gamma &= t
  \end{aligned}
  \right.
\]
matrice
\[
  \begin{pmatrix}
    2 & b & b' \\
    b & b^2 - 2c & t'' \\
    b' & t'' & b'^2 - 2c'
  \end{pmatrix}
\]

\page{184}
avec \ill{}
\[
  (1 \wedge \text{\struck{$x$}}\, x \wedge y \wedge z) \otimes e^{-1} = 2\varepsilon_0
  \qquad \text{où} \quad
  \varepsilon_0 = (\alpha \wedge \alpha' \wedge x_0 \wedge y_0) \otimes e^{-1} \in \pm 1 ,
\]
le signe dépend de la \uncertain{convention} d'orientation (qu'il faudra préciser le moment
venu !).

D'autre part
\[
  \tau - \tau' = \mathrm{Tr}\, \underbrace{[x,y]}_{z} . z = \mathrm{Tr}\, z^2 = \mathrm{Tr}\, 1 = 2
\]
\note{sous $\tau - \tau'$ une première écriture biffée : \struck{$\mathrm{Tr}\, xy \ldots = 1$}.}
d'où
\[
  \Delta = 2\varepsilon\varepsilon_0 \quad \text{i.e.} \quad \varepsilon = \tfrac12\, \varepsilon_0
\]
d'où
\[
  \boxed{\tau - \tau' = \tfrac12\, \varepsilon_0\, \underbrace{(1 \wedge x \wedge y \wedge z) \otimes e^{-1}}_{\overset{\text{déf}}{=}\ \Delta}}
  \qquad \text{i.e.} \quad 2(\tau - \tau') = \varepsilon_0 \Delta
\]
\note{devant $\varepsilon_0$, dans les deux formules, le coefficient $\tfrac12$ est surchargé sur un autre chiffre ($\tfrac14$ ?) ; la ligne « d'où » porte $\Delta$ écrit sur un $2$. La p.~188 écrira $\tau - \tau' = \varepsilon_0 \Delta$.}

(a un sens si $2$ inv.\ dans l'anneau de base $k$ où on travaille).

On en tire
\[
  4\tau\tau' = (\tau+\tau')^2 - (\tau-\tau')^2 = \Sigma^2 - \tfrac14 \Delta^2
\]
\note{ici encore le $\tfrac14$ est surchargé.}
\marginal{?}
À vérifier (dans $\mathbb{Q}[[x_{ij}, y_{ij}, z_{ij}]]$, $1 \leq i,j \leq 2$)\note{dans les crochets, avant $x_{ij}$, une première écriture biffée : \struck{$b, b', b''$}.}
\[
  \Sigma^2 - \tfrac14 \Delta^2 \in 4\, \mathbb{Z}[b,b',b'',t,t',t'']
\]
\struck{i.e.}
\[
  \text{\struck{$4\Sigma^2 - \Delta^2 \in 16\, \mathbb{Z}[b,b',b'',t,t',t'']$}}
\]
d'où l'expression
\[
  \tau\tau' = \Pi(b,b',b'',t,t',t'') \in \mathbb{Z}[b,b',b'',t,t',t'']
\]
expression polynomiale de poids total $6$ (les $b$ étant de
\note{la phrase reprend en tête de la p.~186 ; la p.~185 est un calcul intercalé.}

\page{185}
On va plutôt essayer des calculs
\[
  \tau' - \tau = \mathrm{Tr}\, xyz - \mathrm{Tr}\, yxz = \mathrm{Tr}([x,y]z)
  = \mathrm{Tr}([y,z]x) = \mathrm{Tr}([z,x].y)
\]
(c'est une expression alternée en $x, y, z$, \add{multilinéaire}), donc
\[
  \tau' - \tau = \text{\struck{$\ill{}$}}\, (a \wedge x \wedge y \wedge z) \otimes e^{-1}
\]
pour $a \in \mathcal{A}$ bien déterminé. Comme cette expression est invariante par
$G = \underline{\mathrm{Aut}}\, \mathcal{A}$, \struck{on voit que} et que (si $g \in G(-)$)
\[
  g(a) \wedge g(x) \wedge g(y) \wedge g(z) = a \wedge x \wedge y \wedge z
  \quad (= a \wedge g(x) \wedge g(y) \wedge g(z)) ,
\]
on en conclut que $a = g(a)$ i.e.\ $a$ inv.\ sous $G$ i.e.\ $a = \varepsilon . 1$.
On \ill{} donc \uncertain{aussi}
\[
  \tau' - \tau = \varepsilon\, (1 \wedge x \wedge y \wedge z) \otimes e^{-1}
\]
et pour calculer $\varepsilon$, il suffit de prendre $x, y, z$ tels que
$1 \wedge x \wedge y \wedge z$ \uncertain{inv.}, i.e.\ $1, x, y, z$ une base
\add{de $\mathcal{A}$}. On prendra
\[
  x_0 = \begin{pmatrix} 0 & 0 \\ 1 & 0 \end{pmatrix}, \quad
  y_0 = \begin{pmatrix} 0 & 1 \\ 0 & 0 \end{pmatrix}, \quad
  z_0 = [x_0, y_0] = \begin{pmatrix} -1 & 0 \\ 0 & 1 \end{pmatrix} = 1 - 2\alpha
  \quad \text{où } \alpha = \begin{pmatrix} 1 & 0 \\ 0 & 0 \end{pmatrix}
\]
(NB
$x_0 y_0 = \begin{pmatrix} 0 & 0 \\ 0 & 1 \end{pmatrix}$,
$y_0 x_0 = \begin{pmatrix} 1 & 0 \\ 0 & 0 \end{pmatrix}$)
\note{à la suite, une égalité biffée : $x_0 y_0 z_0 = \begin{pmatrix} 0 & 0 \\ 0 & 1 \end{pmatrix}$.}
On trouve
\[
  1 \wedge x_0 \wedge y_0 \wedge z_0 = -2\, 1 \wedge x_0 \wedge y_0 \wedge \alpha
  = -2\, \alpha' \wedge x_0 \wedge y_0 \wedge \alpha
  = 2\, \alpha \wedge \alpha' \wedge x_0 \wedge y_0
\]
\note{après le dernier $2$, une première écriture biffée : \struck{$x \wedge y \wedge \alpha \wedge \alpha'$}.}
où
\[
  \alpha = \begin{pmatrix} 1 & 0 \\ 0 & 0 \end{pmatrix}, \quad
  \alpha' = \begin{pmatrix} 0 & 0 \\ 0 & 1 \end{pmatrix}
\]
et \struck{$x, y$,} $\alpha, \alpha', x_0, y_0$ forment base can.\ de $M_2(\mathbb{Z})$. On
\note{la suite (« on a donc aussi bien ») est en tête de la p.~184 : la feuille 185 s'insère logiquement avant la p.~184.}

\page{186}
poids $1$ et les $t$ de poids $2$), et de poids $2$ en chacune des
\struck{paires de} variables $x, y, z$ (i.e.\ \ill{} \add{de poids $1$ dans $b, t', t''$,
poids $0$ dans $b', b'', t$}) \struck{$(b,t)\ (b',t')\ (b'',t'')$.} De plus, il doit
être symétrique par permutation des \struck{\ill{}} \add{trois} paires de variables
$(b,t), (b',t'), (b'',t'')$. Et il en va de même pour $\Delta^2$.

On doit donc avoir
\[
  \begin{aligned}
    \Delta^2 = \alpha\, tt't''
    &+ \beta\,(b^2t^2 + b'^2t'^2 + b''^2t''^2)
    + \gamma\,(b'b''t't'' + b''bt''t + bb'tt') \\
    &+ \delta\,(b'b''b^2 t + b''bb'^2 t' + bb'b''^2 t'')
    + \eta\, b^2 b'^2 b''^2
  \end{aligned}
\]
\note{au-dessous, une première version de cette formule, biffée de deux traits et en partie illisible.}
avec
\[
  \alpha, \beta, \gamma, \delta \in \mathbb{Z}
\]
à déterminer, par (\uncertain{quatre} ?) choix particuliers des triples $(x,y,z)$.
P.~ex.\ le triple particulier $x_0, y_0, z_0 = [x_0, y_0]$ donne
\[
  b, b', b'' = 0, \quad \text{et} \quad \Delta = 2\varepsilon_0, \quad \Delta^2 = 4 ,
\]
d'où
\[
  4 = \alpha\, tt't''
\]
où
\[
  t'' = \mathrm{Tr}\, x_0 y_0 = 1, \quad t = \mathrm{Tr}\, y_0 z_0 = 0, \quad t' = \mathrm{Tr}\, z_0 x_0 = 0
\]
et on trouve $4 = \alpha . 0$, ce qui est fâcheux !

Et pour cause --- il faudra, dans l'expression de $\Delta^2$, faire intervenir aussi
$c, c', c''$. En fait, on trouve immédiatement, une fois qu'on y pense :

\page{187}
\[
  \Delta(x,y,z)^2 = \det
  \begin{pmatrix}
    \varphi(1,1) & \varphi(1,x) & \varphi(1,y) & \varphi(1,z) \\
    \varphi(x,1) & \varphi(x,x) & \varphi(x,y) & \varphi(x,z) \\
    \varphi(y,1) & \varphi(y,x) & \varphi(y,y) & \varphi(y,z) \\
    \varphi(z,1) & \varphi(z,x) & \varphi(z,y) & \varphi(z,z)
  \end{pmatrix}
\]
où $\varphi(u,v)$ est la forme bil.\ symm.\ associée à la forme quadratique $q$
($= \det$)
\[
  \varphi(u,v) = \mathrm{Tr}\, u\, \mathrm{Tr}\, v - \mathrm{Tr}(uv) =
\]

\emph{NB} \uncertain{Les}
\[
  \begin{aligned}
    \varphi(u,1) &= \mathrm{Tr}(u) \quad \text{\struck{$\ill{}$}} \\
    \varphi(u,u) &= 2 \det u
  \end{aligned}
\]

Donc on trouve
\[
  \boxed{\Delta(x,y,z)^2 = \det
  \begin{pmatrix}
    2 & b & b' & b'' \\
    b & 2c & bb' - t'' & bb'' - t' \\
    b' & bb' - t'' & 2c' & b'b'' - t \\
    b'' & bb'' - t' & b'b'' - t & 2c''
  \end{pmatrix}}
\]
ce qui permet d'expliciter
\[
  4\tau\tau' = \Sigma^2 - \Delta^2
\]
\marginal{$\longleftarrow$ \ill{}}
où $\Delta$ est donné ci-dessus, et
\[
  \boxed{\Sigma = bt + b't' + b''t'' - bb'b''}
\]

\page{188}
Le calcul « bête », en développant le déterminant en tenant compte des symétries
(par regroupement des termes qui se correspondent par \struck{symétrie} op.\ de
$\mathfrak{S}_3$) et $\Sigma^2$,
\[
  \Sigma^2 - \Delta^2 = 4\Pi(b,b',b'',t,t',t'',c,c',c'')
\]
avec
\[
  \boxed{
  \begin{aligned}
    \Pi = tt't'' &+ (ct^2 + c't'^2 + c''t''^2) + (b^2c'c'' + b'^2c''c + b''^2cc') \\
    &- (b'b''ct + b''bc't' + bb'c''t'')
  \end{aligned}}
\]
\[
  = tt't'' + ct(t - b'b'') + c't'(t' - b''b) + c''t''(t'' - bb')
  + (b^2c'c'' + b'^2c''c + b''^2cc')
\]

On aura, \ill{}
\[
  \left\{
  \begin{aligned}
    \tau + \tau' &= \Sigma \quad (= bt + b't' + b''t'' - bb'b'') \\
    \tau\tau' &= \Pi \quad \bigl(= tt't'' + (ct^2 + c't'^2 + c''t''^2) \\
    &\qquad\quad + (b^2c'c'' + b'^2c''c + b''^2cc') \\
    &\qquad\quad - (b'b''ct + b''bc't' + bb'c''t'')\bigr)
  \end{aligned}
  \right.
\]
\[
  \tau - \tau' = \varepsilon_0\, \Delta
\]

\page{189}
\emph{À vérifier} : L'anneau $B$ \struck{des} des fonctions $F(x,y,z)$ en trois
variables dans $\mathcal{A}$, $G$-invariantes \add{sur} ($\mathcal{A}$ alg.\
d'Azumaya de rang $2$ sur $k$), $G = \underline{\mathrm{Aut}}(\mathcal{A}) \simeq
\mathcal{A}^* / \mathbb{G}_m$), est une algèbre \struck{libre} \add{finie libre} de
rang $2$ sur l'anneau de polynômes à neuf variables
\[
  A = k[\underline{b}, \underline{b}', \underline{b}'', \underline{c}, \underline{c}', \underline{c}'',
  \underline{t}, \underline{t}', \underline{t}'']
\]
où
\[
  \left\{
  \begin{aligned}
    &\underline{b} = \mathrm{Tr}\, x, \quad \underline{b}' = \mathrm{Tr}\, y, \quad \underline{b}'' = \mathrm{Tr}\, z \\
    &\underline{c} = \det x, \quad \underline{c}' = \det y, \quad \underline{c}'' = \det \uncertain{y} \\
    &\underline{t} = \mathrm{Tr}\, yz, \quad \underline{t}' = \mathrm{Tr}\, zx, \quad \underline{t}'' = \mathrm{Tr}\, xy
  \end{aligned}
  \right.
\]
\note{pour $\underline{c}''$ la page porte $\det y$ (on attend $\det z$) ; la lecture de cette lettre est douteuse.}
engendrée par la fonction
\[
  \tau = \mathrm{Tr}\, xyz \quad (= \mathrm{Tr}\, yzx = \mathrm{Tr}\, zxy)
\]
soumise à la relation quadratique
\[
  \boxed{\tau^2 - \Sigma\tau + \Pi = 0}
\]
où \struck{$\Sigma, \Pi$} $\Sigma, \Pi \in A$ sont donnés plus haut. Si $u \mapsto u'$
est l'involution canonique de
\[
  B \simeq A[\underline{\tau}] / (\underline{\tau}^2 - \Sigma\underline{\tau} + \Pi) ,
\]

\page{190}
définie par
\[
  \tau \longmapsto \tau' = \Sigma - \tau ,
\]
on aura
\[
  \tau' = \mathrm{Tr}\, yxz = \mathrm{Tr}\, xzy = \mathrm{Tr}\, zyx .
\]
Les éléments du groupe $\mathfrak{S}_3$ échangent entre eux les neuf générateurs
$\underline{b}, \underline{c}, \underline{t}$, et aussi $\tau$ et $\tau'$ (les
permutations circulaires laissent $\tau, \tau'$ fixes, les transpositions échangent
$\tau, \tau'$). En \struck{\ill{}} tout, $\mathfrak{S}_3$ opère sur $B$ comme
\uncertain{algèbre}, on peut décrire $B$ comme
\[
  B = k[\underline{b}, \underline{b}', \underline{b}'', \underline{t}, \underline{t}', \underline{t}'',
  \underline{c}, \underline{c}', \underline{c}'', \underline{\tau}, \underline{\tau}']
  / (\underbrace{\underline{\tau} + \underline{\tau}' - \Sigma}, \underbrace{\tau\tau' - \Pi})
\]
(équations séparément invariantes sous $\mathfrak{S}_3$) avec l'opération évidente de
$\mathfrak{S}_3$ sur l'anneau de polynômes à onze variables $\underline{b},
\underline{c}, \underline{t}, \underline{\tau}, \underline{\tau}'$.

\emph{NB} Quand $2$ est inv.\ dans $k$, et donc en

\page{191}
car seulement, on peut choisir comme générateur de $B$ sur $A$
\[
  \Delta = \tau - \tau'
\]
soumis à la relation
\[
  \Delta^2 = \Sigma^2 - 4\Pi \in A .
\]
On aura
\[
  \begin{aligned}
    \tau &= \tfrac12 (\Delta + \Sigma) \\
    \tau' &= \tfrac12 (-\Delta + \Sigma)
  \end{aligned}
\]

Mais en car.\ $2$, on aura
\[
  \Delta = \tau + \tau' = \Sigma \in A \quad (\text{en car.\ } 2)
\]
donc $\Delta$ ne risque pas d'\uncertain{être} générateur de $B$ sur $A$ en car.\ $2$,
donc pas non plus si $2$ non inv.\ dans $k$ i.e.\ $k/2k \neq 0$.

\section{Cas de quatre variables}
\note{titre donné par la ligne de la p.~191 : « Cas de quatre variables $x_1, y_1, z_1, t_1$ ».}

Cas de quatre variables $x_1, y_1, z_1, t_1$.
\note{les indices sont lus tels qu'écrits ; la suite note les quatre variables $x_1, x_2, x_3, x_4$.}
Soit $B_4(\mathcal{A}) \subset \mathrm{Sym}(\check{\mathcal{A}}^4)$ l'anneau des
fonctions polynômes $F(x,y,z,t)$

\page{192}
invariantes par $G$. On y trouve les éléments
\[
  \left.
  \begin{aligned}
    b_i &= \mathrm{Tr}\, x_i \\
    c_i &= \det x_i
  \end{aligned}
  \right\} \quad 1 \leq i \leq 4
\]
\[
  t_{ij} = \text{\struck{$t$}}\, t_{\{i,j\}} = \mathrm{Tr}\, x_i x_j
  \qquad \{i,j\} \in \mathfrak{P}_2([1,4])
\]
(ce qui fait déjà $4 + 4 + 6 = 14$ fonctions sur un « espace quotient » de dimension
$4 \times 4 - 3 = 13$ --- il doit y avoir \struck{\ill{}} \add{une} relation entre ces
quantités), et en plus les
\[
  t_{ijk} = \mathrm{Tr}(x_i x_j x_k)
  \qquad i, j, k \in [1,4] \text{ deux à deux distincts}
\]
\note{sous « deux à deux distincts », une première condition biffée : \struck{$1 \leq i < j < k \leq 4$}.}
qui sont \uncertain{algébriques} sur l'anneau engendré par les $b_i, c_i, t_{ij}$.
Il y a $24$ fonctions $t_{ijk}$, \struck{\ill{}} correspondant aux copies de $\mathfrak{S}_3$
(\ill{} $4$ éléments) mais
\[
  t_{ijk} = t_{jki} = t_{kij}
\]
donc il n'y a que $8$ fonctions distinctes, en paires se groupant $2$ par $2$ en
\[
  (*) \quad
  \left\{
  \begin{aligned}
    t_{ijk} + t_{jik} &= \Sigma_{\{i,j,k\}} \in \mathbb{Z}[b_i, c_i, t_i, b_j, c_j, t_j, b_k, c_k, t_k] \\
    t_{ijk} . t_{jik} &= \Pi_{\{i,j,k\}} \in \quad \text{id.}
  \end{aligned}
  \right.
\]
\note{le $\mathbb{Z}$ est écrit sur un $k$ ; les indices des $t$ dans les crochets ne sont pas précisés davantage.}

\page{193}
(fonctions \struck{symétriques} \add{\uncertain{déterminées}} par l'action des groupes
$\mathfrak{S}_{\{i,j,k\}} \subset \mathfrak{S}_4$). Si fait --- en fait, il suffit comme
générateurs les quatre quantités (en plus des générateurs $b_\cdot, c_\cdot, t_{\cdot\cdot}$
du début)
\[
  t_{2,3,4}, \quad t_{1,3,4}, \quad t_{1,2,4}, \quad t_{1,2,3} .
\]
Si $A$ est l'anneau engendré par les $b_\cdot, c_\cdot, t_{\cdot\cdot}$, je présume que $B$
est \add{\uncertain{fini libre}} de rang $2^4$ sur $A$, avec $8$ \add{distincts $t_{ijk}$}
générateurs sur $A$ soumis aux seules relations
\note{l'accolade sous « distincts $t_{ijk}$ » renvoie à « $8$ » ; la phrase semble hésiter entre $8$ et $4$ générateurs.}
\uncertain{(*)} ci-dessus.

Mais il faudrait trouver la relation entre les $b_i, c_j, t_{h,\ell}$ --- peut-être
via les $t_{ijk}$. La clef pour la trouver \add{sera} le calcul de l'expression
\[
  \mathrm{Tr}\, x_1 x_2 x_3 x_4
\]

\page{194}
\note{les calculs de cette page sont écrits sur le verso d'une feuille dactylographiée ; plusieurs blocs sont barrés de longs traits obliques, sans qu'on puisse dire s'il s'agit d'annulations.}
\[
  \beta' = -\beta
\]
\[
  \begin{aligned}
    &bc' + b'c'' + b''c \\
    &- (bc'' + b'c + b''c')
  \end{aligned}
\]
\[
  \alpha' = \qquad \overline{\alpha}' = \overline{\alpha} + \beta\sigma
\]
\[
  \mathbb{Z} \qquad \tau\tau' = \qquad (\alpha - \overline{\alpha})' = \alpha - \overline{\alpha}
\]
\[
  \begin{matrix} b, c, t \\ b', c', t' \\ b'', c'', t'' \end{matrix}
  \qquad \alpha' - \alpha = \beta\sigma
\]
\note{les accents sur $\alpha$ sont tantôt un prime, tantôt une barre ; lecture incertaine.}

Dans une région délimitée par une courbe :
\[
  A \quad \mathfrak{S}_3 \text{ y opère fidèlement}
\]
\[
  b \mapsto b' \mapsto b'' \quad c \mapsto c' \mapsto c'' \quad t \mapsto t' \mapsto t''
\]
\[
  \text{\struck{$\ill{}$}}\ \underline{\mathbb{A}}^3 \times \underline{\mathbb{A}}^3 \times \underline{\mathbb{A}}^3
\]
Fonctions $G$-invariantes en $3$ variables $x, y, z$.

\begin{tikzcd}[column sep=small, row sep=small, nodes={font=\scriptsize}]
  B & & \\
  A \arrow[u, no head] \arrow[r, no head, "\mathfrak{S}_2"] \arrow[d, no head, "\mathfrak{S}_3"'] & B \arrow[d, no head, "\mathfrak{S}_3"] & \\
  A^{\mathfrak{S}_3} \arrow[r, no head, "\mathfrak{S}_2"'] \arrow[d, no head] & B_0 \arrow[d, no head] & \\
  \mathbb{Z}[\Sigma, \Pi] \arrow[r, no head] & B_{00} &
\end{tikzcd}

\note{diagramme d'inclusions, sans flèches ; le $\mathbb{Z}$ du coin inférieur est écrit sur un $k$ biffé. À droite :}
\[
  \mathfrak{S}_3 \times \mathfrak{S}_2 \text{ opère sur } B \qquad \text{i.e.\ } A^{\mathfrak{S}_3}
\]
\[
  \mathfrak{S}_3 \longrightarrow \mathfrak{S}_3 \times \mathfrak{S}_2 \qquad (\mathrm{id}, \mathrm{sg})
\]

\page{196}
\note{verso d'une feuille dactylographiée, comme p.~194.}
\[
  \mathrm{Tr}\, u^n = P_n(\mathrm{Tr}\, u, \det u)
\]
\[
  \mathrm{Tr}\, u^{n_1} v^{m_1} u^{n_2} v^{m_2} \cdots u^{n_r} v^{m_r}
  = P_{(n_1, m_1, n_2, m_2, \ldots, n_r, m_r)}(\ill{})
\]
\[
  \lambda^n + \mu^n = P_n(\lambda + \mu, \lambda\mu)
\]
\[
  \mathrm{Tr}\, x_1 x_2 x_3 x_4 = \mathrm{Tr}\, x_1 x_2 (x_3 x_4)
\]
\[
  \text{\struck{$t_{1234}^2 - \Sigma_{1,2,(3,4)}\, t_{1234} - \cdots$}}
\]
\[
  \tau^2 - \Sigma_{1,2,(3,4)}\, \tau + \Pi_{1,2,(3,4)} = 0
\]
\[
  \Sigma_{1,2,(3,4)}, \ \Pi_{1,2,(3,4)} \in
  \text{\struck{$\mathbb{Z}[b_1, c_1, b_2, c_2$}} \ \text{\add{exprimés en} \ill{}}
\]
\uncertain{fonctions} en
\[
  \underbrace{\mathrm{Tr}\, x_1}_{b_1}, \ \underbrace{\det x_1}_{c_1}, \
  \underbrace{\mathrm{Tr}\, x_2}_{b_2}, \ \underbrace{\det x_2}_{c_2}, \
  \underbrace{\mathrm{Tr}\, x_3 x_4}_{t_{3,4}}, \ \underbrace{\det(x_3 x_4)}_{c_3 c_4}
\]
\[
  \underbrace{\mathrm{Tr}\, x_1 x_2}_{t_{1,2}}, \quad
  \underbrace{\mathrm{Tr}\, x_2(x_3 x_4)}_{t_{2,3,4}}, \quad
  \underbrace{\mathrm{Tr}(x_3 x_4) x_1}_{t_{1,3,4}}
\]
donc dans
\[
  \mathbb{Z}[b_1, c_1, b_2, c_2, c_3, c_4, t_{1,2}, \underbrace{t_{2,3,4}, t_{1,3,4}}_{9}]
\]
\uncertain{entier} sur
\[
  \mathbb{Z}[b_1, c_1, b_2, c_2, b_3, c_3, b_4, c_4, t_{1,2}, t_{1,3}, t_{1,4},
  t_{2,3}, t_{2,4}, t_{3,4}]
\]
\note{le dernier indice est écrit « $t_{24}$ » deux fois ; on attend $t_{3,4}$.}
$\tau$ de degré au plus huit sur $A$ (« \uncertain{divisant} »).

\page{198}
\note{feuille de brouillon, couverte de calculs dispersés et de croquis ; on transcrit les formules dans l'ordre de lecture, de haut en bas.}
\[
  \tau = \mathrm{Tr}(x_1 x_2) x_3 x_4 \qquad x_2 x_1
\]
\[
  \tau^2 - \Sigma_{(1,2),3,4}\, \tau + \Pi
\]
\[
  \mathrm{Tr}\, x_1 x_2 x_3 x_4 + \mathrm{Tr}\, x_2 x_1 x_3 x_4
\]
\[
  x_1 x_2 x_3 x_4 \quad x_2 x_3 x_4 x_1 \quad x_3 x_4 x_1 x_2 \quad x_4 x_1 x_2 x_3
\]
\[
  x_4 x_3 x_2 x_1 \quad x_3 x_2 x_1 \ldots \qquad \tau
\]
\drawing{à droite, un parallélépipède en perspective : sommets marqués par des cercles, certaines arêtes renforcées à l'encre brune, d'autres en pointillé, flèches de rotation autour de deux axes ; étiquettes « $t\ldots$ », « $c$ ». À côté : « $\Sigma_a, \Pi_a$ ». Au-dessus, un petit losange et une flèche circulaire.}
\[
  \underbrace{b_s, c_s, t_s} \ \big| \ \underbrace{t_a} \ \big| \ t_f
\]
\marginal{NB $b_s = b_{s'}$, $c_s = c_{s'}$, $t_a = t_{a'}$, mais $t_s \neq t_{s'}$, $t_f \neq t_{f'}$ ;
de degré $2$ sur $A^{\mathfrak{S}_4}$, de degré $8$ sur $A^{\mathfrak{S}_4}$, de degré $6$ sur
$B^{\mathfrak{S}_4}$ --- de degré $8$ (?) sur $A$}
\note{marge droite très serrée ; les exposants $\mathfrak{S}_4$ sont douteux, et l'ordre des lignes est douteux.}
\[
  \tau^2 \qquad \tau^8 + \alpha_1 \tau^7 + \cdots + \alpha_8 = 0
  \qquad \alpha_i \in \mathbb{Z}[b_\cdot, c_\cdot, t_{\cdot\cdot}]
\]
\[
  g\tau, \ g'\tau, \ g''\tau \qquad \textstyle\sum g\tau \qquad \sum g
\]
\[
  \tau^8 + \beta_1 \tau^7 + \cdots + \beta_8 = 0
\]
\[
  \underbrace{(\alpha_1 - \beta_1)}_{\gamma_1} \tau^7 + \cdots + \underbrace{\alpha_8 - \beta_8} = 0
\]
\drawing{un petit carré aux sommets numérotés $1, 2, 3, 4$ ; à côté : « $F(x, \underline{y, z}, t)$ ».}
\[
  H \subset \mathfrak{S}_4 \qquad B^H \overset{\text{degré } 6}{\text{---}} B^{\mathfrak{S}_4} = B^{\mathfrak{S}}
\]
\[
  \alpha + \beta\tau \quad \text{à } b \text{ stables par } \mathfrak{S}_3^+
\]
\[
  \alpha' + \beta'\tau' = \alpha' + \beta'(\sigma - \tau) = (\alpha' + \beta'\sigma) - \beta'\tau
\]
\[
  \beta' = -\beta \ / \ \alpha' + \beta\sigma = \alpha
  \qquad \alpha' = \alpha + \beta\sigma \ / \ (\alpha')' = \alpha = \alpha' \ldots
\]

\page{199}
\note{verso d'une feuille dactylographiée (« L 22 »), écrite de travers ; on ne transcrit que l'encre de G.}

\begin{tikzcd}[column sep=small, row sep=small, nodes={font=\scriptsize}]
  A \arrow[r] & B & \\
  A^{\mathfrak{S}_3^+} \arrow[u, "3"] \arrow[r, no head, "2"] & B_1 = B^{\mathfrak{S}_3^+} \arrow[u, "3"'] & \\
  A^{\mathfrak{S}_3} \arrow[u, "2"] \arrow[r, no head, "2"'] & B_0 \arrow[u, "2"'] & B^{\mathfrak{S}_3}
\end{tikzcd}

\note{$B^{\mathfrak{S}_3}$ est placé entre $A^{\mathfrak{S}_3}$ et $B_1$, relié aux deux par des traits ; une forme en pointillé au crayon encadre le bas du diagramme.}
\[
  \begin{aligned}
    \alpha + \beta\tau + \alpha' + \beta'\tau'
    &= (\alpha + \alpha') + \bigl(\beta\tau + \beta'(\sigma - \tau)\bigr) \\
    &= \underbrace{\alpha + \alpha'} + \underbrace{\beta'\sigma} + (\beta - \beta')\tau
  \end{aligned}
\]
\[
  \alpha + \alpha' + \beta\sigma - (\alpha + \alpha' + \beta'\sigma) = (\beta - \beta')\sigma
  \qquad \text{OK}
\]

\page{200}
\note{verso d'une feuille dactylographiée, comme les pp.~194--199.}
\[
  \begin{aligned}
    \tau_f + \tau_{f'} &= \Sigma_a \\
    \tau_{f'} + \tau_{f''} &= \Sigma_{a'}
  \end{aligned}
  \qquad \longrightarrow \qquad
  \tau_f - \tau_{f''} = \Sigma_a - \Sigma_{a'} = \Sigma_{\ill{}} - \Sigma_{\ill{}}
\]
\[
  \text{i.e.} \quad \Sigma_{\ill{}} + \Sigma_{\ill{}} = \Sigma_{a'} + \Sigma_{\ill{}}
\]
\drawing{cube en perspective, avec arêtes renforcées ($a$, $a'$), sommets cerclés, flèches de rotation et étiquettes $f$, $f'$, $f''$.}
\[
  \Sigma_a + \Sigma_{\ill{}} = \Sigma_{a'} + \Sigma_{\ill{}}
\]
\[
  \tau_f + \tau_{f'} = \Sigma \qquad \tau_{\tilde f} + \tau_{\tilde f'} = \Sigma_{\ill{}} - \Sigma
\]
\[
  4 . 3 . 2 = 24 = 48/2 \qquad
  \tau_f \tau_{f'} = \Pi \qquad \tau_{\tilde f} \tau_{\tilde f'} = \Pi \ldots
\]
\[
  \text{donc } \{\tau_f, \tau_{f'}\} = \{\tau_{\tilde f}, \tau_{\tilde f'}\} .
\]
\drawing{second cube, arêtes $b$, $f'$, $t''$ marquées.}
\[
  \tau_f - \tau_{\ill{}} = \Sigma \ldots = \Sigma_b - \Sigma_{b'}, \qquad
  \text{donc } \tau_f \neq \tau_{\tilde f} \ \text{(\uncertain{or})}, \quad \tau_f = \tau_{\tilde f'}
\]
\[
  \Sigma_{\ill{}} + \Sigma_b = \Sigma_{b'} + \Sigma_{\ill{}}
\]
\[
  \begin{aligned}
    \Sigma_{\ill{}} - \Sigma_{\ill{}} &= \Sigma_{\ill{}} - \Sigma_{\ill{}} \\
    \Sigma_{\ill{}} + \Sigma_{\ill{}} &= \Sigma_{\ill{}} - \Sigma_{\ill{}}
  \end{aligned}
\]
$\Sigma_{\ill{}}$ ne dépend que de la diagonale de \ill{}, donc il n'y a que
trois \emph{valeurs} de $\Sigma_a$ ; donc il n'y a que trois \uncertain{valeurs} de $\Pi_a$.
\drawing{troisième cube, arêtes étiquetées $\Sigma_1$, $\Sigma_2$, $\Sigma_3$.}

$12$ relations de la forme
\[
  \Sigma_a + \Sigma_{a'} = \Sigma_b + \Sigma_{b'}
\]
\struck{pour $\Sigma_a$ fixé} via opérations du groupe symétrique $\mathfrak{S}_4$, elles
se \uncertain{déduisent} \ill{} de l'une \ill{}
\note{la phrase se poursuit au-delà de la fin du lot.}

\end{document}
